% =====================================================================
%  Issue register — Collagen / Gallic-Acid Crosslinking (Graph_model)
%
%  Every entry was reproduced by executing the released code on the
%  released data. Entries marked [pre-existing] were additionally
%  reproduced on pristine upstream/main @ 5cb24be with no local edits.
%
%  Compile: pdflatex issues_by_category.tex
%  Core LaTeX only - no siunitx / tcolorbox / mhchem.
% =====================================================================
\documentclass[11pt,a4paper]{article}

\usepackage[margin=2.3cm]{geometry}
\usepackage[T1]{fontenc}
\usepackage{booktabs}
\usepackage{longtable}
\usepackage{array}
\newcolumntype{L}[1]{>{\raggedright\arraybackslash}p{#1}}
\usepackage{xcolor}
\usepackage{amsmath}
\usepackage[colorlinks=true,linkcolor=blue,urlcolor=blue]{hyperref}

\definecolor{okgreen}{RGB}{0,120,60}
\definecolor{badred}{RGB}{175,0,30}
\definecolor{warnamber}{RGB}{180,110,0}

\newcommand{\ok}[1]{\textcolor{okgreen}{#1}}
\newcommand{\bad}[1]{\textcolor{badred}{#1}}
\newcommand{\warn}[1]{\textcolor{warnamber}{#1}}
\newcommand{\code}[1]{\texttt{\small\hyphenchar\font=`\-\relax #1}}

% severity tags
\newcommand{\SBLOCK}{\bad{\textbf{[BLOCKING]}}}
\newcommand{\SMAJOR}{\warn{\textbf{[MAJOR]}}}
\newcommand{\SMINOR}{\textbf{[MINOR]}}
\newcommand{\SFIXED}{\ok{\textbf{[FIXED]}}}
\newcommand{\SOPEN}{\bad{\textbf{[OPEN]}}}

\title{Issue Register\\
       \large Graph\_model --- Programming/Model defects and Biochemical data errors}
\author{Reproducibility audit, executed against released code and data}
\date{2026-08-17}

\begin{document}
\sloppy
\maketitle

\section*{How to read this}

Issues are split into two registers:

\begin{itemize}\itemsep2pt
  \item \textbf{Section 1 --- Programming and Model.} Software defects and
        experimental-design problems. These affect whether a number can be
        produced or trusted, but do not misdescribe chemistry.
  \item \textbf{Section 2 --- False biochemical data.} Entries where the
        stored data does not represent the intended molecule, protein, or
        physical process. These are wrong \emph{as chemistry}, independent of
        any code.
\end{itemize}

Status tags: \SFIXED{} corrected in this branch; \SOPEN{} still live.
Severity: \SBLOCK{} invalidates results; \SMAJOR{} changes numbers;
\SMINOR{} correctness/hygiene. \textbf{[pre-existing]} means it was
reproduced on pristine \code{upstream/main} and was not introduced by
this audit.

\vspace{4pt}\noindent
\textbf{Headline:} under leave-one-ligand-out cross-validation on the
released data, Model~A reaches Spearman $\rho = 0.163$ against a reported
$0.93$, with RMSE $1.047$ against a global-mean baseline of $0.992$ and a
per-docking-box baseline of $0.918$.

% =====================================================================
\section{Programming and Model}
% =====================================================================

\subsection{P1. Defects that made results impossible or false}

\begin{longtable}{@{}L{0.6cm}L{5.2cm}L{8.6cm}@{}}
\toprule
\textbf{\#} & \textbf{Issue} & \textbf{Detail and consequence} \\
\midrule
\endhead
P1.1 & \SBLOCK{} \SFIXED{} Silent NaN training &
  \code{train\_lolo\_cv} wrapped train, validation \emph{and} test batches in
  \code{except Exception: continue} with no logging. A runtime error killed
  every batch; \code{regression\_metrics([],[])} returned NaN; all 9 folds
  ``completed''. A full LOLO run finished in ${\sim}50$\,s reporting
  \code{rmse\_mean: nan}. [pre-existing] \\
P1.2 & \SBLOCK{} \SFIXED{} Every forward pass of Models A--D raised &
  \code{option\_a.\_encode\_raw} imported \code{ConditionEncoder} from a wrong
  path \emph{inside} \code{forward()}, so it failed at runtime, not import
  time. B, C and D reuse \code{\_encode\_raw}. Combined with P1.1 this is why
  no result could be produced. [pre-existing] \\
P1.3 & \SBLOCK{} \SFIXED{} Package unimportable &
  The \code{features/} directory sat at \code{Graph\_model/features/} while
  every import expected \code{Graph\_model/data/features/}. \code{import
  Graph\_model.data.config} raised \code{ModuleNotFoundError}. Restored by
  upstream \code{69a6c14}; the imports were always right, the tree was wrong. \\
P1.4 & \SMAJOR{} \SFIXED{} \code{biopython} absent from requirements &
  Required by \code{level2\_protein.py}. Without it the builder catches the
  \code{ImportError} and substitutes a \emph{single all-zero residue node}.
  A by-the-book install gives every sample an empty receptor. Verified:
  \code{residue.x} is \code{[1,30]} zeros without it, \code{[30,30]} with. \\
P1.5 & \SMAJOR{} \SOPEN{} Published metrics mix two target quantities &
  Training logs record \code{n\_train}\,$=9517$, \code{n\_val}\,$=1679$
  (total $11{,}196 = 6{,}196$ anchor $+\,5{,}000$ PDBbind; \code{--max-transfer}
  defaults to 5000). About 45\,\% of the validation set is PDBbind, so the
  reported kcal/mol RMSE averages Vinardo docking scores with PDBbind
  experimental affinities. \\
P1.6 & \SMAJOR{} \SOPEN{} Documented loss never executed &
  All four training loops use \code{nn.MSELoss()}.
  \code{CombinedDockingLoss} (MSE\,+\,ListMLE\,+\,monotonicity\,+\,NLL) is
  defined and exported but \emph{instantiated nowhere}, so the ranking and
  monotonicity terms and the heteroscedastic uncertainty head are untrained.
  The main-text ``smooth-L1'' description matches neither. \\
P1.7 & \SMAJOR{} \SFIXED{} Results overwritten in place &
  \code{train\_main.py} rewrites \code{Graph\_model/results/} with no warning
  and no merge. A single-model run replaces the A--I comparison file,
  deleting the other eight entries and the Table~IV provenance block. The
  command in the project README has this property. \\
P1.8 & \SMAJOR{} \SOPEN{} Shipped cache is stale &
  The v2 bundle's \code{.pkl} is byte-identical to v1
  (md5 \code{a6e4b35d...}). Since \code{use\_cache=True} is the default,
  training silently uses the \emph{old} 40-record MMP-1 data unless
  \code{--force-reload} is passed. \\
\bottomrule
\end{longtable}

\subsection{P2. Reproducibility and environment}

\begin{longtable}{@{}L{0.6cm}L{5.2cm}L{8.6cm}@{}}
\toprule
\textbf{\#} & \textbf{Issue} & \textbf{Detail and consequence} \\
\midrule
\endhead
P2.1 & \SMAJOR{} \SFIXED{} Hardcoded absolute data root &
  \code{REPO\_ROOT = Path("/Users/suppboat/Jupyter\_Dock")}, and
  \code{PROCESSED\_DIR.mkdir()} executed at \emph{import} time, creating
  directories on any machine that imported the package. Now
  \code{\$COLLAGEN\_DATA\_ROOT} with a repo-relative fallback. \\
P2.2 & \SMINOR{} \SFIXED{} Circular import &
  The cycle was \code{features} $\to$ \code{conditions} $\to$
  \code{data.config} $\to$ \code{data/\_\_init\_\_} $\to$ \code{dataset}
  $\to$ \code{features}. Import success depended on which package was
  imported first. Fixed with a
  lazy (PEP 562) package \code{\_\_init\_\_}. \\
P2.3 & \SMAJOR{} \SFIXED{} Incomplete seeding &
  Only \code{torch.manual\_seed}; \code{random}, \code{numpy},
  \code{cuda.manual\_seed\_all} and cuDNN determinism were unseeded, so two
  runs with the same \code{--seed} were not the same run. \\
P2.4 & \SMAJOR{} \SFIXED{} Five disagreeing device resolvers &
  \code{run\_model\_d\_interpret.py} had \emph{no CUDA branch} and silently
  used CPU on NVIDIA hardware; \code{hpo.py} cleared the allocator cache only
  on CUDA, so MPS sweeps leaked memory; \code{train\_main.resolve\_device}
  accepted any string, so \code{--device cuda} on a Mac failed later with an
  unrelated error. Unified in \code{Graph\_model/train/device.py}. \\
P2.5 & \SMINOR{} \SOPEN{} Hardware claim &
  SI states an H100 GPU; the released training logs record
  \code{device: cpu}. \\
P2.6 & \SMAJOR{} \SFIXED{} No entry point, no tests &
  Zero \code{\_\_main\_\_} guards, zero \code{argparse}, a 1-byte README, and
  five 1--5 byte placeholder files named \code{test}. \code{pytest} appeared
  zero times. Upstream later added \code{train\_main.py}; this branch adds
  \code{scripts/run\_lolo.py} and a 70-test suite. \\
P2.7 & \SMINOR{} \SOPEN{} Dependency pinning &
  All requirements are \code{>=} with no lockfile. \\
\bottomrule
\end{longtable}

\subsection{P3. Data-handling defects}

\begin{longtable}{@{}L{0.6cm}L{5.2cm}L{8.6cm}@{}}
\toprule
\textbf{\#} & \textbf{Issue} & \textbf{Detail and consequence} \\
\midrule
\endhead
P3.1 & \SMAJOR{} \SFIXED{} Missing targets imputed as \code{0.0} &
  Against a $-3.83 \pm 1.07$ kcal/mol distribution this is a fabricated
  $+3.5\sigma$ outlier. \ok{Did not fire on the released data (0 imputed
  values), so no published number is affected.} Now drops the row. \\
P3.2 & \SMAJOR{} \SFIXED{} Silent row loss &
  \code{hetero\_dataset} counted build failures behind five \code{debug}
  messages. Rows could vanish without trace. Now logs every failure with its
  reason and refuses to build if $>5\,\%$ of rows drop. \\
P3.3 & \SMINOR{} \SOPEN{} Dead dataset path &
  \code{CollagenDockingDataset} emits a homogeneous 54-D \code{Data} object;
  every model indexes \code{data['ligand']}, which it never sets. It cannot
  feed any model. \\
P3.4 & \SMINOR{} \SFIXED{} Feature-dimension confusion &
  Five different values appeared for ``the node feature dimension'': 35, 54,
  74, 79 and a docstring claiming 32/33. Only two featurisers exist:
  \textbf{35}-D (\code{graph/level1\_ligand.py}, consumed by every model) and
  54-D (\code{data/features/atom.py}, dead path). Now asserted at runtime. \\
\bottomrule
\end{longtable}

\subsection{P4. Model and experimental design}

\begin{longtable}{@{}L{0.6cm}L{5.2cm}L{8.6cm}@{}}
\toprule
\textbf{\#} & \textbf{Issue} & \textbf{Detail and consequence} \\
\midrule
\endhead
P4.1 & \SBLOCK{} \SOPEN{} Rank correlation does not reproduce &
  Converged LOLO-CV, anchor-only, seed 0: Spearman $\rho = 0.163 \pm 0.074$
  against Table~IV's $0.93$. Folds early-stopped at epochs 2--73 (200-epoch
  cap never reached), and raising the budget from 10 to 200 epochs changed
  $\rho$ by $-0.004$, so undertraining does not explain it. \\
P4.2 & \SBLOCK{} \SOPEN{} No baselines reported &
  Measured LOLO RMSE $1.047$ versus a global-mean baseline of $0.992$ and a
  per-docking-box mean of $0.918$. Only 4 of 9 folds beat their own baseline.
  A model that cannot beat a per-box mean has not demonstrated learned
  chemistry. \\
P4.3 & \SMAJOR{} \SOPEN{} Signature of overfitting &
  Training longer made generalisation \emph{worse}: LOLO RMSE rose from
  $0.938$ (10-epoch cap) to $1.047$ (converged) while $\rho$ stayed flat.
  Fold validation RMSE (held-in ligands) sits at 0.52--0.70 against test RMSE
  (held-out ligand) of 0.40--2.37. \\
P4.4 & \SMAJOR{} \SOPEN{} Eight of nine models never read the protein &
  Only Option~B references \code{data['residue']} / \code{interacts}. A, C,
  D, E, F, G, H, I consume atom features, bond edges and a condition vector
  only. They are ligand-only regressors over nine unique molecules. \\
P4.5 & \SMAJOR{} \SOPEN{} Models F and G have no real geometry &
  EGNN coordinates are \code{Linear(35$\to$3)} of atom features; DimeNet
  ``angles'' come from \code{cosine\_similarity} of edge features and
  ``distances'' from \code{norm(edge\_attr)}. Genuine RDKit ETKDG conformers
  exist in \code{features/conformer\_3d.py} and are imported nowhere. Neither
  model is SE(3)-equivariant as built. \\
P4.6 & \SMAJOR{} \SOPEN{} Leaky split reported &
  \code{StratifiedSplitter} (row-level 70/15/15) is imported by nothing, yet
  the ``best epoch'' and validation-RMSE columns are row-level artifacts.
  With nine molecules a row-level split leaks molecular identity --- as
  \code{lolo\_cv.py} states in its own docstring. \\
P4.7 & \SMAJOR{} \SOPEN{} Model D checkpoint misread &
  $\text{val\_RMSE}\approx 0.549$ is a single-fold row-level value. The
  per-ligand mean baseline on the same data is $0.589$, so this figure is
  consistent with memorising each ligand's mean. \\
P4.8 & \SMINOR{} \SOPEN{} Headline ligand is the worst fold &
  Pentagalloylglucose: test RMSE $2.370$, $\rho = 0.083$ --- more than twice
  the next-worst fold. A single averaged RMSE conceals this. \\
P4.9 & \SMINOR{} \SOPEN{} No $R^2$, no reference model &
  \code{metrics.py} computes RMSE/MAE/Pearson/Spearman only. \code{r2\_score},
  \code{sklearn} and \code{xgboost} appear zero times despite being declared
  dependencies. No baseline of any kind is implemented. \\
P4.10 & \SMINOR{} \SOPEN{} Unexercised modules &
  \code{maml.py}, \code{contrastive.py} and \code{hpo.py} are reached by no
  entry point. If Optuna HPO selected on the validation split, the reported
  validation metrics are optimistically biased. \\
\bottomrule
\end{longtable}

% =====================================================================
\section{False biochemical data}
% =====================================================================

These entries are wrong as \emph{chemistry or structural biology}: the
stored value does not represent the intended molecule, protein, or physical
process. All were verified with RDKit or Bio.PDB against the released files.

\subsection{B1. Ligand structures}

\begin{longtable}{@{}L{0.6cm}L{4.0cm}L{4.6cm}L{5.1cm}@{}}
\toprule
\textbf{\#} & \textbf{Entry} & \textbf{As released} & \textbf{Correct} \\
\midrule
\endhead
B1.1 & \SBLOCK{} \SFIXED{} Pentagalloylglucose &
  \bad{6 disconnected fragments}; 114 heavy atoms; MW 1609.15;
  $\mathrm{C_{68}H_{56}O_{46}}$. A \code{.}-separated string of five free
  gallic acids plus a partially galloylated ring. &
  \ok{1 fragment}; 67 heavy atoms; MW 940.68;
  $\mathrm{C_{41}H_{32}O_{26}}$. The headline ligand was fed to every model
  as six separate molecules, ${\sim}1.7\times$ too large. \\
B1.2 & \SMAJOR{} \SFIXED{} \code{PGG\_SMILES\_RDKIT} &
  A cyclopentane pentagallate: no pyranose ring oxygen, no C6
  hydroxymethyl. $\mathrm{C_{40}H_{30}O_{25}}$, MW 910.66, 65 atoms. &
  Connected \emph{but still the wrong molecule}. Read by nothing
  (\code{smiles\_rdkit} had zero references), so adopting it as the ``fix''
  would have passed a fragment-count test while remaining wrong. \\
B1.3 & \SMAJOR{} \SFIXED{} Ellagic acid &
  \bad{Monolactone}: $\mathrm{C_{13}H_{8}O_{6}}$, 19 heavy atoms, MW 260.20
  --- missing one of the two lactone bridges. &
  \ok{Dilactone}: $\mathrm{C_{14}H_{6}O_{8}}$, 22 heavy atoms, MW 302.19.
  The entry's own comment said ``dilactone'' and its \code{mw}/\code{n\_ha}
  described the correct molecule. \\
B1.4 & \SMINOR{} \SFIXED{} EDC &
  MW 191.70 --- the \emph{hydrochloride salt} --- while the SMILES is the
  free base. & 155.24 ($\mathrm{C_8H_{17}N_3}$), matching the SMILES. \\
B1.5 & \SMINOR{} \SFIXED{} EDC O-acylisourea &
  \code{n\_ha} 13. & 16. (Also a propionic-acid stand-in, not the gallic-acid
  adduct.) \\
B1.6 & \SMAJOR{} \SFIXED{} NHS ester intermediate &
  MW 285.22 and \code{n\_ha} 10 --- matching neither each other nor the
  SMILES ($\mathrm{C_{10}H_{16}N_2O_4}$, MW 228.25, 16 atoms). &
  Resolved by reading the shipped docking inputs. The EDC/NHS series is built
  on \emph{propionic acid} as the carboxylic-acid stand-in (the catalogue and
  the structure file agree on the O-propanoylisourea), so the consistent
  NHS-stage species is \ok{N-succinimidyl propionate},
  $\mathrm{C_7H_9NO_4}$, MW 171.15, 12 atoms. Catalogue now
  \ok{9/9 self-consistent}. \\
\bottomrule
\end{longtable}

\subsection{B2. Chemical feature detection}

\begin{longtable}{@{}L{0.6cm}L{5.2cm}L{8.6cm}@{}}
\toprule
\textbf{\#} & \textbf{Issue} & \textbf{Detail} \\
\midrule
\endhead
B2.1 & \SMAJOR{} \SFIXED{} Galloyl SMARTS matched nothing &
  \code{\_SMARTS\_GALLOYL = "Oc1cc(O)c(O)cc1"} encodes
  1,3,4-trihydroxybenzene (hydroxyquinol), not a 3,4,5-galloyl ring: the
  hydroxyls are not three-adjacent. \code{galloyl\_strict} is \bad{0 for all
  nine ligands, gallic acid included}, so the feature the mechanism rests on
  never fires. Correct pattern:
  \code{[OX2H]c1cc(cc([OX2H])c1[OX2H])[CX3]=[OX1]} --- a trihydroxyphenyl
  ring bearing an acyl carbon, i.e.\ 3,4,5-trihydroxybenzoyl. Now
  \ok{gallic acid 1, PGG 5, pyrogallol 0} (no acyl group, so not a galloyl
  unit). The strict \code{xfail} was removed after it XPASSed. \\
B2.2 & \SMINOR{} \SFIXED{} Galloyl weight double-counted &
  Because B2.1 never fired, \code{galloyl\_weighted} was carried entirely by
  the catechol and pyrogallol patterns, which \emph{both} match each galloyl
  ring: $1.00 + 0.67 = 1.67$ per ring (gallic acid $\to 1.67$, PGG
  $\to 8.35$). The root cause was counting per SMARTS \emph{hit} and
  subtracting (\code{ca - g - py}); a benzene-1,2,3-triol yields \emph{two}
  catechol hits, so one spurious catechol survived on every trihydroxy ring.
  Counting is now done at \ok{ring level with disjoint classes}: gallic acid
  1.00, pyrogallol 1.00, protocatechuic acid 0.67, ellagic acid 1.34,
  PGG 5.00 --- all matching reference chemistry. \\
B2.4 & \SMAJOR{} \SFIXED{} Phenolic OH over-counted &
  \code{total\_aromatic\_oh} accepted \emph{any} oxygen of degree 1--2 bonded
  to an aromatic carbon, with no hydrogen requirement. RDKit perceives
  ellagic acid's fused lactone rings as aromatic, so its two ester oxygens
  and two lactone carbonyls were scored as phenolic hydroxyls: \bad{8
  against 4 actual}. Now requires a hydroxyl hydrogen. All nine ligands match
  reference values (gallic acid 3, pyrogallol 3, protocatechuic acid 2,
  ellagic acid 4, PGG 15). \\
B2.3 & \SMINOR{} \SFIXED{} Fragment count disagreed with catalogue &
  With the released PGG SMILES the detector found \bad{9} galloyl units
  against a declared 5, because the five free gallic acids were counted
  separately. Now 5. \\
\bottomrule
\end{longtable}

\subsection{B3. Protonation physics}

\begin{longtable}{@{}L{0.6cm}L{5.2cm}L{8.6cm}@{}}
\toprule
\textbf{\#} & \textbf{Issue} & \textbf{Detail} \\
\midrule
\endhead
B3.1 & \SMAJOR{} \SFIXED{} pH encoding violated Henderson--Hasselbalch &
  \code{PROPKA\_PROTONATION} hardcodes $f = \{5.0{:}\,0.85,\ 5.5{:}\,0.15,\
  7.0{:}\,0.02\}$. No single p$K_a$ produces an $0.85 \to 0.15$ drop across
  0.5 pH units: fitting $f=0.85$ at pH 5.0 implies p$K_a = 5.75$, which
  predicts $0.64$ at pH 5.5 and $0.05$ at pH 7.0. The drop is
  thermodynamically impossible for one titratable site --- over half a pH
  unit, Henderson--Hasselbalch bounds any single site to at most
  $0.76 \to 0.24$, and only when centred exactly on its p$K_a$.
  The fractions are now \ok{derived} from one p$K_a$ via
  $f = 1/(1 + 10^{\,\mathrm{pH} - \mathrm{p}K_a})$, giving
  $\{5.0{:}\,0.1051,\ 5.5{:}\,0.0358,\ 7.0{:}\,0.0012\}$, with
  $\mathrm{p}K_a = 4.07$ for GLU (Olsson \emph{et al.}, \emph{JCTC} 7(2):525,
  2011 --- the reference already used elsewhere in the codebase). A test
  asserts the table equals the formula, so a hand-typed value cannot
  return. \\
B3.2 & \SMAJOR{} \SFIXED{} Two contradictory pH models &
  \code{graph/residue\_data.py} already implemented Henderson--Hasselbalch
  correctly, but the \emph{incorrect} table was the one reaching the network.
  There is now \ok{one pH model}: \code{config.protonation\_fraction} and
  \code{residue\_data.protonation\_fraction} are asserted equal for GLU at
  every pH, and \code{GLU\_PKA} is tied to the same referenced p$K_a$ table.
  Off-study pH values now evaluate the titration curve instead of snapping to
  the nearest tabulated point, which had made the feature a step function
  where the physics is smooth. \emph{Caveat recorded in the code}: the
  corrected fractions span 0.001--0.105, ${\sim}8\times$ narrower than the
  old (wrong) 0.02--0.85, so the encoding is now right but poorly conditioned
  as a network input; the principled reparameterisation is the HH exponent
  $(\mathrm{p}K_a - \mathrm{pH})$, which is a \emph{modelling} change and was
  therefore not applied. \\
\bottomrule
\end{longtable}

\subsection{B4. Receptor structures and docking boxes}

\begin{longtable}{@{}L{0.6cm}L{5.2cm}L{8.6cm}@{}}
\toprule
\textbf{\#} & \textbf{Issue} & \textbf{Detail} \\
\midrule
\endhead
B4.1 & \SBLOCK{} \SOPEN{} Every MMP-1 record has an empty receptor &
  \bad{40/40} (v1) and \bad{120/120} (v2) MMP-1 graphs have \emph{zero}
  residue nodes. Reproduced identically on pristine \code{upstream/main}, and
  present in the authors' own shipped cache. [pre-existing] \\
B4.2 & \SBLOCK{} \SOPEN{} Root cause: missing box geometry &
  The MMP-1 CSV has no \code{box\_center\_x/y/z} or \code{box\_size\_A}
  columns (the collagen CSV does). \code{hetero\_dataset.py:222} falls back
  to a 20\,\AA{} box at the \emph{origin}; the MMP-1 protein occupies
  $z \in [13.45,\,57.92]$, so the box never intersects the structure.
  Reconstructed true centres lie a mean of 65.5\,\AA{} from the origin. \\
B4.3 & \SMAJOR{} \SOPEN{} Box labels are lossy &
  The pose tree contains \textbf{67} distinct MMP-1 boxes collapsed to 8
  canonical labels. \code{mmp1\_box\_mapping\_to\_canonical8.json} records
  remapping distances of \bad{20.79, 32.77 and 39.43\,\AA}: a row labelled
  \code{ASP\_GLU\_cluster1} was mapped to \code{LYS\_cluster42}, a different
  site 20.8\,\AA{} away. \\
B4.4 & \SMAJOR{} \SOPEN{} Receptor identity is unrecorded &
  Poses exist against \code{966C} (296 res), \code{966C\_apo} (157 res) and
  \code{1SU3} (1225 res, a different crystal in a different coordinate
  frame). \code{level2\_protein.py} always selects \code{966C}, so
  \code{1SU3} coordinates would be interpreted in the wrong frame. No CSV
  column identifies which receptor a row used. \\
B4.5 & \SMAJOR{} \SOPEN{} MMP-1 coverage cannot support a selectivity claim &
  5 of 9 ligands (no EDC/NHS species), a \emph{single} pH (5.5) and a
  \emph{single} temperature (25\,\textdegree C) --- so the pH and temperature
  inputs are constant and carry no information for MMP-1. Under LOLO-CV,
  4 of 9 folds contain zero MMP-1 test records. \\
B4.6 & \SMINOR{} \SOPEN{} Documented MMP-1 counts are wrong &
  \code{README} justifies Model~D's $\times 10$ loss weight with ``120 MMP-1
  vs $\sim$2{,}052 collagen''. Measured: 40 (v1) or 120 (v2) versus
  \textbf{6{,}156} collagen --- a ratio of $1{:}154$ or $1{:}51$, not
  $1{:}17$. \\
B4.7 & \SMINOR{} \SOPEN{} Collagen binding sites are pre-cropped &
  \code{pig\_collagen\_I\_alpha2\_\allowbreak pH5.5\_H\_fix\_\allowbreak bindingsite\_clean.pdb}
  contains \textbf{3 residues / 41 atoms}. Collagen graphs therefore carry
  only 3 or 6 residue nodes, so Option~B's cross-attention attends over 3--6
  residues in total. \\
\bottomrule
\end{longtable}

\subsection{B6. Ligand identity: the docked molecule is not always the named one}

\emph{Found 2026-08-17.} Nothing in the pipeline ever compared the molecule
named in the catalogue against the molecule the docking campaign actually read
from disk. Reading
\code{<data\_root>/Phukhao/collagen\_gallic\_results/*.sdf} and the 335
released poses shows \bad{three of nine ligands were docked as a different
compound}. Identical in the v1 and v2 Drive snapshots, so this is not a
regression introduced by the dataset update.

\begin{longtable}{@{}L{0.6cm}L{3.4cm}L{5.0cm}L{4.9cm}@{}}
\toprule
\textbf{\#} & \textbf{Ligand} & \textbf{Named (catalogue)} & \textbf{Actually docked} \\
\midrule
\endhead
B6.1 & \SBLOCK{} \SOPEN{} Pentagalloylglucose &
  $\mathrm{C_{41}H_{32}O_{26}}$, MW 940.68, 67 heavy atoms,
  \textbf{5 galloyl units}. &
  The input \code{.sdf} holds \bad{4 disconnected fragments}
  ($\mathrm{C_{48}H_{42}O_{33}}$): a tri-\emph{O}-galloylglucose plus three
  loose gallic acids. Docking kept only the largest, so every released pose is
  \bad{tri-\emph{O}-galloylglucose}, $\mathrm{C_{27}H_{24}O_{18}}$, MW 636.47,
  45 heavy atoms, \textbf{3 galloyl units}. Confirmed on all 68 PGG poses. \\
B6.2 & \SMAJOR{} \SOPEN{} NHS &
  N-hydroxysuccinimide, $\mathrm{C_4H_5NO_3}$, MW 115.09 (CID 6467). &
  $\mathrm{C_3H_3NO_3}$, MW 101.06 --- \bad{one $\mathrm{CH_2}$ short}. \\
B6.3 & \SMAJOR{} \SOPEN{} NHS ester intermediate &
  N-succinimidyl propionate, $\mathrm{C_7H_9NO_4}$, MW 171.15, five-membered
  imide ring. &
  $\mathrm{C_6H_7NO_4}$, MW 157.12, with a \bad{four-membered ring} ---
  again exactly one $\mathrm{CH_2}$ short. \\
\bottomrule
\end{longtable}

B6.2 and B6.3 are \emph{one} defect, not two: a succinimide ring has four ring
carbons and both structures were built with three. The deficiency is exactly
$\mathrm{CH_2}$ (14.03\,Da) in each case, so a single upstream
structure-building error propagated to every NHS-containing species.

\paragraph{Consequence.} A binding affinity is only meaningful beside the
structure it was computed for. Every $\Delta G$ labelled
\emph{pentagalloylglucose} belongs to a molecule carrying three of the five
galloyl groups whose count is the paper's proposed mechanism --- and PGG is the
headline ligand. The remaining six ligands (gallic acid, EDC, EDC
\emph{O}-acylisourea, protocatechuic acid, pyrogallol, ellagic acid)
\ok{match exactly} and are unaffected.

\paragraph{Remedy.} Re-dock the three structures. The scores cannot be
corrected by editing the catalogue, and were left untouched. To stop this
recurring, \code{Graph\_model/data/provenance.py} now compares every catalogue
entry against its docking input by molecular formula and InChIKey skeleton;
it runs in the training preflight, is recorded in each results file under
\code{provenance.ligand\_identity}, and is asserted by
\code{tests/test\_provenance.py}. A mismatch outside the three known ones
fails the suite.

\subsection{B5. Recoverability}

Box centres for MMP-1 \emph{can} be reconstructed from the 335 released
docked poses: a pose can only be generated inside its search box, so the poses
are physical evidence of where the box was.
\code{scripts/reconstruct\_mmp1\_boxes.py} implements this and was run.

\begin{itemize}\itemsep1pt
  \item \textbf{Validated on collagen}, where the true centres are recorded:
        31 boxes, \ok{mean error 4.94\,\AA}, median 4.78\,\AA{}, range
        1.39--7.94\,\AA{}. Against a 20\,\AA{} box this recovers the right
        binding \emph{region}, not the original box.
  \item \textbf{Coverage}: \ok{67 of 67} (receptor, box) pairs in the CSV
        recovered --- 45 for \code{porcine\_MMP1\_1SU3}, 8 for
        \code{porcine\_MMP1\_966C}, 14 for \code{porcine\_MMP1\_966C\_apo}.
  \item \textbf{Occupancy}, the test that matters. The defect is not that the
        centre is missing but that the resulting box is empty. At the origin
        (current behaviour) all three receptors give \bad{0 atoms}. The
        reconstructed centres give a median of \ok{148--220 receptor atoms}
        per box, with \ok{61 of 67} boxes non-empty.
  \item \textbf{Six boxes remain unusable} --- \code{1SU3/GLU\_cluster1},
        \code{1SU3/GLU\_cluster15}, and four on \code{966C\_apo}. Their poses
        sit off the receptor surface, so the centroid is not a box centre.
        These need the original configuration files.
  \item \textbf{Not recoverable at all}: the box \emph{size}. Pose extent is a
        lower bound, not a measurement, and is emitted under a name that says
        so.
\end{itemize}

Output goes to a \emph{sidecar} CSV carrying \code{method} and
\code{is\_measured=false} columns. The shipped CSVs are not modified and the
reconstruction is \textbf{not} wired into training: reconstructed geometry is
not measured geometry, and whether to use it is the authors' call. The
authoritative values should still come from the original Vina/Vinardo
configuration files, which additionally resolve B4.3 (which box) and B4.4
(which receptor). Synthetic or placeholder coordinates must not be
substituted: an invented centre does not leave the receptor absent, it puts
the model on the \emph{wrong} residues while appearing to work.

% =====================================================================
\section{Priority}
% =====================================================================

\begin{enumerate}\itemsep2pt
  \item \textbf{P4.1, P4.2} --- reconcile the reported $\rho = 0.93$ with the
        measured $0.163$, and publish both baselines. Nothing else in the
        manuscript matters until this is settled.
  \item \textbf{P1.5} --- state whether Table~IV is anchor-only; if not, the
        kcal/mol figures are not collagen $\Delta G$ metrics.
  \item \textbf{B6.1--B6.3} --- re-dock pentagalloylglucose, NHS and the NHS
        ester intermediate. Until then no affinity for those three ligands can
        be reported under those names, and the galloyl-count mechanism has no
        support from the PGG data.
  \item \textbf{B3.1, B2.1} --- \emph{done}: the pH table and the galloyl
        SMARTS are both corrected. Every feature vector changed, so all
        feature-dependent numbers need one re-run under feature schema~2.
  \item \textbf{B4.2} --- obtain the Vina configs; until then no MMP-1
        structural claim is supportable. Reconstructed centres (B5) cover 61
        of 67 boxes at ${\sim}5$\,\AA{} and are available as a disclosed
        sidecar in the meantime.
  \item \textbf{P1.6, P4.4, P4.5} --- correct the manuscript's descriptions of
        the loss, protein usage, and SE(3) equivariance to match the code.
\end{enumerate}

\end{document}
